The ablations qualify a simple claim that either temporal decay or
encoder richness is sufficient to explain the gains. A constant
article kernel already reduces QLIKE by 7.16\%; learned decay
improves further, with linear decay producing the lowest losses in
the kernel comparison. Sentiment probabilities combined with learned
decay nearly match the full-dimensional CryptoBERT reference, and
the 256-component PCA representation improves on it in this run.
These results support compact text representations and explicit
temporal weighting as useful design choices, while showing that
the exponential, full-dimensional specification is not uniquely
favored by the forecast comparison.
Forecast improvements also persist with lookbacks of six, twelve,
and twenty-four hours, with six hours giving the lowest mean losses
in the lookback comparison.
